% ═══════════════════════════════════════════════════════════════
% SLIDES: Las preposiciones de lugar
% Spanisch Klasse 8 - Unidad 6, DS 2 (45 Min)
% ═══════════════════════════════════════════════════════════════

\documentclass[aspectratio=169,11pt]{beamer}

\usepackage[utf8]{inputenc}
\usepackage[T1]{fontenc}
\usepackage[ngerman]{babel}
\usepackage{helvet}
\usepackage{tikz}
\usepackage{tabularx}
\usepackage{booktabs}
\usepackage{amssymb}

% ═══════════════════════════════════════════════════════════════
% THEME & FARBEN
% ═══════════════════════════════════════════════════════════════

\usetheme{default}
\usecolortheme{default}
\setbeamertemplate{navigation symbols}{}
\setbeamertemplate{footline}[frame number]

\definecolor{spanisch}{HTML}{C41E3A}
\definecolor{hellgrau}{HTML}{F5F5F5}
\definecolor{dunkelgrau}{HTML}{333333}

\setbeamercolor{frametitle}{fg=white,bg=spanisch}
\setbeamercolor{title}{fg=white,bg=spanisch}
\setbeamercolor{structure}{fg=spanisch}
\setbeamercolor{block title}{fg=white,bg=spanisch}
\setbeamercolor{block body}{bg=hellgrau}

\setbeamerfont{frametitle}{series=\bfseries}
\setbeamerfont{title}{series=\bfseries}

% Befehle
\newcommand{\es}[1]{\textcolor{spanisch}{\textbf{#1}}}
\newcommand{\de}[1]{\textcolor{dunkelgrau}{\textit{#1}}}
\newcommand{\phase}[1]{\begin{center}\colorbox{spanisch}{\textcolor{white}{\large\bfseries #1}}\end{center}}

% ═══════════════════════════════════════════════════════════════
% DOKUMENT
% ═══════════════════════════════════════════════════════════════

\begin{document}

% ---------------------------------------------------------------
% TITEL
% ---------------------------------------------------------------
\begin{frame}
\begin{center}
\vspace{1cm}
{\Huge\bfseries\textcolor{spanisch}{Las preposiciones de lugar}}\\[0.5cm]
{\Large Die Ortspräpositionen}\\[1cm]
{\normalsize Spanisch Klasse 8 · Unidad 6 · Qué pasa 1, S. 106-107}
\end{center}
\end{frame}

% ---------------------------------------------------------------
% LERNZIELE
% ---------------------------------------------------------------
\begin{frame}{Hoy aprendemos...}
\begin{itemize}
\setlength\itemsep{1em}
\item[\textcolor{spanisch}{\textbullet}] \textbf{Ortspräpositionen verstehen:} delante de, detrás de, entre...
\item[\textcolor{spanisch}{\textbullet}] \textbf{Positionen beschreiben:} ¿Dónde está el perro?
\item[\textcolor{spanisch}{\textbullet}] \textbf{Ein Stadtviertel beschreiben:} El banco está enfrente de la iglesia.
\end{itemize}
\end{frame}

% ═══════════════════════════════════════════════════════════════
% PHASE 1: EINSTIEG (5 Min)
% ═══════════════════════════════════════════════════════════════

\begin{frame}{}
\phase{Einstieg: Wiederholung}
\end{frame}

\begin{frame}{Kurze Wiederholung: Los lugares}
\begin{center}
{\Large Gestern haben wir gelernt:}\\[0.8cm]
\begin{columns}
\column{0.5\textwidth}
\begin{itemize}
\item \es{el parque}
\item \es{la biblioteca}
\item \es{el supermercado}
\item \es{la panadería}
\end{itemize}

\column{0.5\textwidth}
\begin{itemize}
\item \es{el banco}
\item \es{la farmacia}
\item \es{el cine}
\item \es{la iglesia}
\end{itemize}
\end{columns}
\vspace{0.8cm}
{\Large\bfseries Heute: Wo ist was?}
\end{center}
\end{frame}

% ═══════════════════════════════════════════════════════════════
% PHASE 2: INPUT - PRÄPOSITIONEN (10 Min)
% ═══════════════════════════════════════════════════════════════

\begin{frame}{}
\phase{Input: Las preposiciones de lugar}
\end{frame}

\begin{frame}{Die Ortspräpositionen (1/2)}
\begin{columns}
\column{0.5\textwidth}
\begin{itemize}
\setlength\itemsep{0.8em}
\item \es{enfrente de} \de{gegenüber von}
\item \es{delante de} \de{vor}
\item \es{detrás de} \de{hinter}
\item \es{encima de} \de{auf/über}
\item \es{debajo de} \de{unter}
\end{itemize}

\column{0.5\textwidth}
\begin{center}
\begin{tikzpicture}[scale=0.8]
% Tisch
\fill[gray!30] (0,0) rectangle (3,0.3);
\draw[thick] (0,0) rectangle (3,0.3);
\draw[thick] (0.3,0) -- (0.3,-1);
\draw[thick] (2.7,0) -- (2.7,-1);
% Ball oben
\fill[spanisch] (1.5,1) circle (0.3);
\node at (1.5,1.6) {\footnotesize encima};
% Ball unten
\fill[spanisch] (1.5,-0.6) circle (0.3);
\node at (1.5,-1.2) {\footnotesize debajo};
\end{tikzpicture}
\end{center}
\end{columns}
\end{frame}

\begin{frame}{Die Ortspräpositionen (2/2)}
\begin{columns}
\column{0.5\textwidth}
\begin{itemize}
\setlength\itemsep{0.8em}
\item \es{al lado de} \de{neben}
\item \es{a la derecha de} \de{rechts von}
\item \es{a la izquierda de} \de{links von}
\item \es{entre} \de{zwischen}
\item \es{en} \de{in}
\end{itemize}

\column{0.5\textwidth}
\begin{center}
\begin{tikzpicture}[scale=0.8]
% Drei Boxen
\fill[gray!30] (0,0) rectangle (1,1);
\node at (0.5,0.5) {A};
\fill[spanisch!30] (1.5,0) rectangle (2.5,1);
\node at (2,0.5) {B};
\fill[gray!30] (3,0) rectangle (4,1);
\node at (3.5,0.5) {C};
% Labels
\node at (2,-0.5) {\footnotesize \es{entre} A y C};
\end{tikzpicture}
\end{center}
\end{columns}
\end{frame}

\begin{frame}{Achtung: de + el = del}
\begin{center}
{\Large Wenn nach \es{de} der Artikel \es{el} kommt:}\\[1cm]
{\LARGE \es{de} + \es{el} $\rightarrow$ \es{del}}\\[1cm]
\pause
{\Large Beispiel:}\\[0.5cm]
{\Large \es{El banco está enfrente \textbf{del} supermercado.}}\\
{\normalsize \de{Die Bank ist gegenüber vom Supermarkt.}}
\end{center}
\end{frame}

\begin{frame}{}
\begin{center}
{\Large\bfseries $\rightarrow$ Lernblatt: Präpositionentabelle ausfüllen}\\[1cm]
{\normalsize Übertrage die deutschen Übersetzungen ins Lernblatt.}
\end{center}
\end{frame}

% ═══════════════════════════════════════════════════════════════
% PHASE 3: ÜBUNG 1 - COPITO (15 Min)
% ═══════════════════════════════════════════════════════════════

\begin{frame}{}
\phase{Übung 1: ¿Dónde está Copito?}
\end{frame}

\begin{frame}{Copito – der kleine Hund}
\begin{center}
{\Large Öffnet das Buch auf Seite 106, Übung 2B}\\[1cm]
{\LARGE \es{¿Dónde está Copito?}}\\
{\normalsize \de{Wo ist Copito (der Hund)?}}\\[1cm]
{\Large Schaut euch die 10 Bilder an.}\\[0.5cm]
{\Large Beschreibt, wo der Hund ist!}
\end{center}
\end{frame}

\begin{frame}{Beispiel: Bild 1}
\begin{center}
\begin{tikzpicture}[scale=1.2]
% Tisch
\fill[gray!30] (0,0) rectangle (3,0.3);
\draw[thick] (0,0) rectangle (3,0.3);
\draw[thick] (0.3,0) -- (0.3,-1.5);
\draw[thick] (2.7,0) -- (2.7,-1.5);
% Hund unter dem Tisch
\fill[spanisch] (1.5,-0.8) circle (0.35);
\node[white] at (1.5,-0.8) {\tiny Copito};
% Label
\node at (1.5,-2) {la mesa};
\end{tikzpicture}

\vspace{0.8cm}
\pause
{\Large\bfseries Antwort:}\\[0.3cm]
{\Large \es{Copito está debajo de la mesa.}}
\end{center}
\end{frame}

\begin{frame}{Jetzt ihr! Bilder 2-10}
\begin{center}
{\Large\bfseries $\rightarrow$ AB Aufgabe 1}\\[0.8cm]
{\Large Schreibt für jedes Bild einen Satz:}\\[0.5cm]
\begin{tabular}{ll}
Bild 2: & Copito está... \\[0.4em]
Bild 3: & Copito está... \\[0.4em]
... & ... \\[0.4em]
Bild 10: & Copito está... \\
\end{tabular}
\\[0.8cm]
{\normalsize Ihr habt 8 Minuten Zeit.}
\end{center}
\end{frame}

\begin{frame}{Lösungen besprechen}
\begin{center}
{\Large Wer möchte vorlesen?}\\[1cm]
\begin{tabular}{rl}
Bild 1: & \es{debajo de la mesa} \\[0.4em]
Bild 2: & \es{encima de la silla} \\[0.4em]
Bild 3: & \es{delante del sofá} \\[0.4em]
Bild 4: & ... \\
\end{tabular}
\end{center}
\end{frame}

% ═══════════════════════════════════════════════════════════════
% PHASE 4: ÜBUNG 2 - MI BARRIO (12 Min)
% ═══════════════════════════════════════════════════════════════

\begin{frame}{}
\phase{Übung 2: Mi barrio}
\end{frame}

\begin{frame}{Höraufgabe: ¿Qué barrio?}
\begin{center}
{\Large Öffnet das Buch auf Seite 107, Übung 4}\\[1cm]
{\LARGE\bfseries Schritt 1: Hören}\\[0.5cm]
{\Large Hört das Audio.}\\[0.3cm]
{\Large Welches Stadtviertel wird beschrieben?}\\[1cm]
\begin{tabular}{cc}
$\square$ \textbf{Viertel A} & $\square$ \textbf{Viertel B} \\
\end{tabular}
\end{center}
\end{frame}

\begin{frame}{Das ANDERE Viertel beschreiben}
\begin{center}
{\LARGE\bfseries Schritt 2: Schreiben}\\[0.5cm]
{\Large\bfseries $\rightarrow$ AB Aufgabe 2}\\[0.8cm]
{\Large Beschreibt das ANDERE Viertel}\\
{\normalsize (das, das NICHT im Audio war)}\\[0.8cm]
{\Large Mindestens 5 Sätze!}\\[0.5cm]
\begin{block}{Redemittel}
\es{El/La \_\_\_ está enfrente de...}\\
\es{El/La \_\_\_ está al lado de...}\\
\es{El/La \_\_\_ está entre... y...}
\end{block}
\end{center}
\end{frame}

\begin{frame}{Beispielsätze}
\begin{center}
{\Large So könnt ihr schreiben:}\\[0.8cm]
\begin{itemize}
\setlength\itemsep{0.5em}
\item \es{El banco está enfrente de la iglesia.}
\item \es{La farmacia está al lado del supermercado.}
\item \es{El parque está entre la biblioteca y el cine.}
\item \es{La panadería está a la derecha del restaurante.}
\end{itemize}
\end{center}
\end{frame}

\begin{frame}{Vorlesen}
\begin{center}
{\Large Wer möchte seine Beschreibung vorlesen?}\\[1cm]
{\normalsize Achtet auf:}\\[0.5cm]
\begin{itemize}
\item Richtige Präposition?
\item \es{de + el = del}?
\item Artikel richtig (\es{el/la})?
\end{itemize}
\end{center}
\end{frame}

% ═══════════════════════════════════════════════════════════════
% PHASE 5: SICHERUNG (3 Min)
% ═══════════════════════════════════════════════════════════════

\begin{frame}{}
\phase{Sicherung}
\end{frame}

\begin{frame}{Was haben wir heute gelernt?}
\begin{itemize}
\setlength\itemsep{1em}
\item[\textcolor{spanisch}{\checkmark}] \textbf{Präpositionen:} enfrente de, delante de, detrás de, al lado de...
\item[\textcolor{spanisch}{\checkmark}] \textbf{Struktur:} ... está + Präposition + Ort
\item[\textcolor{spanisch}{\checkmark}] \textbf{Kontraktion:} de + el = del
\end{itemize}

\vspace{1cm}

\begin{block}{Hausaufgabe}
Beschreibe dein eigenes Viertel (5-6 Sätze) mit den Präpositionen.
\end{block}
\end{frame}

\begin{frame}{}
\begin{center}
{\Huge\bfseries\textcolor{spanisch}{¡Hasta la próxima!}}\\[1cm]
{\Large Bis zum nächsten Mal!}
\end{center}
\end{frame}

\end{document}
